%%%%%%%%%%%%%%%%%%%%%%%%%%%%%%%%%%%%%%%%%%%%%%%%%%
%%%%%%%%%%%%%%%%%%%%%%%%%%%%%%%%%%%%%%%%%%%%%%%%%%
\section{Discussion}\label{remdo:sec:discussion}

\ifdefined\DISSERTATION
    Beyond the analysis of individual results in \Cref{remdo:sec:results}, this section discusses the study's contributions, implications, and limitations more broadly.
    It compares results to existing studies for context and suggests avenues for future work to resolve model shortcomings and answer additional design questions.
\else
    Beyond the per-result analysis of \Cref{remdo:sec:results}, this section discusses the study's contributions and limitations, comparing results to existing studies and suggesting future work.
\fi

%%%%%%%%%%%%%%%%%%%%%%%%%%%%%%%%%%%%%%%%%%%%%%%%%%
\subsection{Result Comparison to Other Studies }
% (and comparison to commercial WECs)
This subsection compares the key findings to prior studies and highlights where the present results align or differ.

%%%%%%%%%%%%%%%%%%%%%%%
\subsubsection{Force Saturation}
\ifdefined\DISSERTATION
    Previous studies \cite{mcgilton_optimal_2024,coe_maybe_2021,devin_high-dimensional_2024,gaebele_tpl_2025,mccabe_force-limited_2024} recommend capping the maximum force well below the unsaturated point, finding that this sacrifices only minimal power, but none perform a full economic analysis confirming that the cost savings outweigh the power loss.
    The results in \Cref{remdo:sec:results-single} verify the conclusion with an appropriate cost analysis, since the optimal design sizes its powertrain small enough to require force saturation.
    The follow-up analysis in \Cref{remdo:sec:sensitivity-model-assumptions} repeats the optimization without allowing force saturation and finds a 5\% LCOE penalty.
    Force saturation is therefore highly valuable as a control technique, capable of tangible LCOE reductions once the power--cost tradeoff is considered.
    Future studies could quantify the lowest marginal PTO force price (\$/N) for which this conclusion still holds.
\else
    Previous studies \cite{mcgilton_optimal_2024,coe_maybe_2021,devin_high-dimensional_2024,gaebele_tpl_2025,mccabe_force-limited_2024} recommend capping maximum force well below the unsaturated point to minimize power loss, but none have economically confirmed that the financial savings outweigh the performance reduction.
    This study verifies the conclusion with a full cost analysis: the optimal design sizes its powertrain to require force saturation, and re-optimizing without it incurs a 5\% LCOE penalty. %(\Cref{remdo:sec:sensitivity-model-assumptions}), confirming force saturation as a valuable control technique.
\fi
% \hl{add similar analysis for power saturation}

%%%%%%%%%%%%%%%%%%%%%%%
\subsubsection{Global Deployability}
\ifdefined\DISSERTATION
    A survey indicates that 40\% of WEC developers design for site-agnostic deployment, so global deployability is of industrial interest \cite{trueworthy_wave_2020}.
    A previous study \cite{de_andres_adaptability_2015} compares the global deployability of a uniform versus a tunable WEC design across wave climates, finding that tuning the resonant period to the dominant ocean period by changing geometry increases capture width ratio but decreases the more important power-to-mass ratio (an inverse LCOE proxy), and concluding that tuning is unfavorable because it requires expensive larger diameters.
    Importantly, however, those devices have equal diameter and height, use only damping control, and are tuned to maximize power without regard for mass or cost.
    In MDOcean, tuning can be accomplished by varying diameter and height independently and with reactive control, and devices are tuned to minimize LCOE rather than maximize power.
    The single-objective results in \Cref{remdo:tab:opt-dv-values,remdo:tab:opt-output} agree with the overall conclusion of \cite{de_andres_adaptability_2015} that power maximization does not align with LCOE minimization, though here the additional cost of maximum power appears as PTO cost rather than structural cost.
    Meanwhile, the location sensitivity (\Cref{remdo:tab:location}) shows a significantly larger optimal float in Hawaii, the site with the lowest common wave period, which departs from the study \cite{de_andres_adaptability_2015}, which finds small (4~m diameter) WECs optimize their LCOE proxy regardless of dominant period.
    This demonstrates that a more thorough multidisciplinary model is essential to draw accurate design conclusions.
\else
    A survey finds 40\% of WEC developers design for site-agnostic deployment \cite{trueworthy_wave_2020}.
    A prior study \cite{de_andres_adaptability_2015} concludes that geometric tuning to the dominant wave period is unfavorable, but its devices have coupled diameter and height, use only damping control, and tune for power rather than cost.
    MDOcean's results agree that power maximization does not align with LCOE minimization, but show the cost penalty as PTO rather than structural cost, and show a significantly larger optimal float in Hawaii (\Cref{remdo:tab:opt-dv-values,remdo:tab:opt-output,remdo:tab:location}), which contradicts \cite{de_andres_adaptability_2015} and demonstrates the value of a fuller multidisciplinary model.
\fi

%%%%%%%%%%%%%%%%%%%%%%%
\subsubsection{Structural Efficiency}
A previous study compares structural mass per unit hull volume for marine structures including WECs and ship hulls, finding $278.76$~kg/m\textsuperscript{3} for steel structures \cite{roberts_bringing_2021}, with a 95\% confidence interval of $\pm280$ metric tons.
\ifdefined\DISSERTATION
    This trendline predicts a nominal float structural mass of 72-632 metric tons and a minimum-LCOE structural mass of 1814-2374 metric tons, whereas MDOcean calculates 213 and 80 metric tons respectively.
    The nominal design lies within the confidence interval, but the optimized design lies far below it.
    Even though the optimal design's hull volume exceeds that of any data point used in the fit, the large departure warrants concern: the only data points with structural mass comparable to the optimized 80-ton float have an order-of-magnitude lower hull volume, indicating that MDOcean may vastly underestimate the required thickness.
    This requires further investigation and may stem from underestimating loads in the dynamics module or overestimating the factor of safety in the structures module.
\else
    This trendline predicts a minimum-LCOE structural mass of 1814-2374 metric tons, whereas MDOcean calculates 80 metric tons, far below the confidence interval.
    Since the only comparable-mass data points have an order-of-magnitude lower hull volume, MDOcean may underestimate the required thickness, warranting further investigation of possible load underestimation or factor-of-safety overestimation.
\fi

% compare to the Sandia RM3 optimization results

%%%%%%%%%%%%%%%%%%%%%%%
\subsubsection{Optimization Procedure}
\label{remdo:sec:discussion-optimization}
\ifdefined\DISSERTATION
    Recalling the prior WEC optimization studies of \Cref{remdo:sec:lit}, all prior large-scale ($\geq 8$ design variables) WEC optimizations use the genetic algorithm, a population-based heuristic.
    Based on the stated generation counts and population sizes \cite{khanal_multi-objective_2024,garcia-teruel_reliability-based_2021,cotten_multi-objective_2022,abdulkadir_control_2024}, these require roughly 12{,}000, 2{,}200, 6{,}100, and 2{,}400 function evaluations to converge.
    The present study is the first large-scale WEC optimization with a gradient-based algorithm, which was expected to require fewer evaluations.
    Generating the Pareto seeds, however, entails 20 single-objective optimizations at about 100 iterations each, using finite difference over 12 design variables, for an estimated $20\cdot100\cdot(12+1)=\resultsRE[paretoFcnEvals]$ evaluations,  dwarfing the genetic algorithm counts.
    This surprising result is attributable to finite difference.
    Despite the greater evaluation count, MDOcean's semi-analytical model likely makes the optimization faster than the BEM-based alternatives, though those papers report no runtime to confirm this.
    Given the fast runtime, little effort was spent reducing iterations, but moderate reductions should be possible by tuning convergence criteria, exploiting a gradient sparsity pattern to cut finite-difference evaluations, and tuning the ratio of seeds to pattern-search set size.
    A more substantial reduction is possible via automatic differentiation or the adjoint method, which would make the evaluation count independent of the number of design variables at the cost of greater implementation effort: the adjoint method would require constraint aggregation given the many constraints, and automatic differentiation would require rewriting functions in terms of elementary operations and manually supplying gradients for unsupported functions such as the Bessel functions in the hydrodynamics module.
    Even with the higher evaluation count, the gradient-based algorithm retains the advantage of certifying convergence to a local optimum, and with systematic multistart, reasonable confidence of a global optimum, which a heuristic optimization cannot provide.
\else
    All prior large-scale ($\geq 8$ design variables) WEC optimizations use the genetic algorithm (\Cref{remdo:sec:lit}), requiring roughly 2{,}000--12{,}000 function evaluations to converge \cite{khanal_multi-objective_2024,garcia-teruel_reliability-based_2021,cotten_multi-objective_2022,abdulkadir_control_2024}.
    This study is the first to use a gradient-based algorithm at this scale; surprisingly, finite-difference gradients over 12 design variables make its estimated $\resultsRE[paretoFcnEvals]$ evaluations exceed the heuristic counts, though MDOcean's semi-analytical model likely still runs faster in wall-clock time.
    Automatic differentiation or the adjoint method would make the count independent of the design-variable dimension, and the gradient-based approach retains the key advantage of certifying local optimality, with global confidence via multistart, that heuristics cannot.
\fi

%%%%%%%%%%%%%%%%%%%%%%%%%%%%%%%%%%%%%%%%%%%%%%%%%%
\subsection{Model Realism}
\ifdefined\DISSERTATION
    Using the optimization and sensitivity results, the modeling assumptions of \Cref{remdo:sec:model} can now be assessed for reasonableness.
    We examine both the model complexities MDOcean captures, to confirm they were worth including, and the shortcomings it cannot capture, to judge whether they must be added in future versions.
    Previous studies such as \cite{garcia-teruel_reliability-based_2021} have called for this kind of systematic model-fidelity assessment to build designer confidence in the most relevant effects and their computational cost.
\else
    The optimization and sensitivity results allow the modeling assumptions from reference \cite{mccabe_development_2026} to be assessed, both the complexities worth including and the shortcomings that may need adding, which is the kind of systematic fidelity assessment called for by the study \cite{garcia-teruel_reliability-based_2021}.
\fi

%%%%%%%%%%%%%%%%%%%%%%%
\subsubsection{Observed Impact of Modeled Effects}
\ifdefined\DISSERTATION
    Several areas of model complexity, force saturation, slamming amplitude limits, and 2-DOF dynamics, stand out as contributing significantly to the results, evident in the constraint activity and in the differences from using 1-DOF (stationary-spar) dynamics; \Cref{remdo:sec:sensitivity-model-assumptions} analyzed these in detail.
    MDOcean's 2-DOF dynamics run roughly an order of magnitude slower than its 1-DOF dynamics \cite{mccabe_development_2026}, suggesting that code improvements to narrow this gap may be worthwhile given the dynamic importance of the second degree of freedom.
    Force saturation is especially relevant: it directly reduces powertrain cost and indirectly reduces structural cost by lowering operational heave loading, which permits a thinner spar given the active column fatigue factor-of-safety constraint.
    The observed impact of these effects justifies the corresponding model complexity and suggests other optimization frameworks should include them for realistic results.
\else
    Force saturation, slamming amplitude limits, and 2-DOF dynamics all contribute significantly to the results.% (\Cref{remdo:sec:sensitivity-model-assumptions}).
    The 2-DOF dynamics run roughly an order of magnitude slower than 1-DOF \cite{mccabe_development_2026}, suggesting worthwhile code improvements, while force saturation reduces both powertrain cost and, via lower heave loading, structural cost.
    The impact of these effects justifies their model complexity and argues for their inclusion in other frameworks.
\fi

\ifdefined\DISSERTATION
    Conversely, a few modeling complexities appear unnecessary in hindsight because they have little effect.
    \Cref{remdo:sec:sensitivities-local} finds low sensitivity to drag coefficient, and the re-optimization without drag in \Cref{remdo:sec:sensitivity-model-assumptions} confirms it, suggesting that modeling the drag nonlinearity is inessential and that future optimizations may disable drag to avoid iterating the drag damping coefficient.
    Likewise, the slamming model is derived to be valid not only for long waves but also intermediate wavelengths, where the maximum amplitude depends on the phase of the dynamic response; the simple long-wave slamming model is reasonably accurate at sea states with an active slamming constraint \cite{mccabe_development_2026}, suggesting the intermediate-wavelength model may be unnecessary.
    Confirming this would require checking all points during the optimization, not just the optimum, since the slamming model fidelity is not currently a model setting and the long-wave condition is verified manually at the optimum. %\hl{I should actually check this.}  % TODO (scrutineering)
    Finally, although drag and intermediate-wavelength slamming each have little individual effect, the problem's nonlinearity admits a possible coupling effect were both disabled simultaneously, or an emergent effect under a changed formulation (a different objective or added constraints), so users should exercise caution when simplifying the model.
\else
    Conversely, 
    %the low sensitivity to drag coefficient (\Cref{remdo:sec:sensitivities-local}), confirmed by re-optimization without drag (\Cref{remdo:sec:sensitivity-model-assumptions}), suggests the drag nonlinearity is inessential.
    %Similarly, 
    the simple long-wave slamming model is reasonably accurate where the slamming constraint is active \cite{mccabe_development_2026}, suggesting the intermediate-wavelength model may be unnecessary.
    Because the problem is nonlinear, however, disabling both simultaneously or changing the formulation could produce coupling effects, so model simplification warrants caution.
\fi

%%%%%%%%%%%%%%%%%%%%%%%
\subsubsection{Potential Impact of Unmodeled Effects}
\ifdefined\DISSERTATION
    No constraints corresponding to priority 1 or 2 requirements are active in the optimal solution, % \hl{(double check this - structures effective breadth could make this false)}  % TODO (scrutineering)
    a hopeful sign that the true optimum is unlikely to lie in a region of model inaccuracy.
    Likewise, the inactivity of the pitch stability constraint mitigates the impact of any metacentric-height error arising from the geometry module's even-mass-distribution assumption for the center of gravity \cite{mccabe_development_2026}.
    On the other hand, the active operational-sea-state slamming constraints mean some assumptions are worth revisiting.
    In particular, MDOcean enforces the float amplitude limits as optimization constraints rather than through the constrained optimal linear controller \cite{mccabe_development_2026}, and assumes the WEC must operate in every sea state with nonzero probability in the joint probability distribution rather than entering a survival mode \cite{mccabe_development_2026}; both assumptions interact with the active amplitude constraints and merit future reconsideration.
\else
    No priority 1 or 2 constraints are active at the optimum, suggesting the true optimum is unlikely to lie in a region of model inaccuracy, and the inactive pitch stability constraint mitigates any metacentric-height error from the geometry module \cite{mccabe_development_2026}.
    However, the active operational slamming constraints invite revisiting two assumptions \cite{mccabe_development_2026}: enforcing amplitude limits as optimization constraints rather than via the constrained optimal controller, and requiring operation in all nonzero-probability sea states rather than allowing a survival mode.
\fi

%%%%%%%%%%%%%%%%%%%%%%%%%%%%%%%%%%%%%%%%%%%%%%%%%%
\subsection{Formulation Realism}
This subsection evaluates how realistic the optimization formulation choices are and what refinements are most impactful.

%%%%%%%%%%%%%%%%%%%%%%%
\subsubsection{Design Variable Suitability}
\ifdefined\DISSERTATION
    \Cref{remdo:sec:formulation} described this study's careful choice of design variables to facilitate satisfaction of high-priority requirements, such as maintaining model validity and aiding convergence, while balancing the computational challenge of a large design space.
    The analysis of \Cref{remdo:sec:sensitivities-local} reveals that several variables chosen as parameters may merit inclusion as design variables in future studies: $D_d$, $T_{f,1}$, $T_s$, and possibly $D_{f,in}$.
    Adding $N_{WEC}$ as a design variable carries little importance in LCOE minimization, where economies of scale always favor larger $N_{WEC}$, but offers useful insight in a modified multi-objective analysis where the first objective changes from per-WEC power to array power.
    This is relevant for small-scale Powering the Blue Economy applications with a fixed power requirement, where it is unclear whether multiple smaller WECs or one larger WEC is preferable.
    
    Additionally, a sensitivity analysis of design variables perturbed as if they were parameters would help quantify inter-disciplinary coupling and potentially inform formulation simplifications.
    For example, a high sensitivity of the optimal structural design variables to the powertrain design variables $\left(\frac{dx^*_{\text{struct}}}{dx_{pto}}\right)$ would support this study's choice of fully coupled MDO \cite{mccabe_development_2026}, while a low value would support disciplinary decoupling.
    Such analysis may be necessary if a future optimization with many more design variables is too costly to perform fully coupled and must be split into sequential optimizations (e.g.\ bulk-geometry power maximization, then powertrain sizing, then structural sizing).
    Performing this sequential optimization for the current formulation and assessing the suboptimality would quantify the value of coupling.
    The authors predict, however, that the speed of the semi-analytical models and the extensions discussed here will permit full coupling even for substantially larger WEC optimizations, including the ambitious target of comparing multiple hydrodynamic architectures with realistic powertrain models.
\else
    \Cref{remdo:sec:formulation} described the careful choice of design variables to maintain model validity and aid convergence.
    The sensitivity analysis (\Cref{remdo:sec:sensitivities-local}) suggests $D_d$, $T_{f,1}$, $T_s$, and possibly $D_{f,in}$ merit inclusion as design variables in future studies.
    Adding $N_{WEC}$ matters little for LCOE minimization but is valuable in a modified multi-objective analysis using array power, relevant to fixed-requirement Powering the Blue Economy applications weighing many small WECs against one large one.
    A sensitivity analysis of design variables perturbed as parameters would further quantify inter-disciplinary coupling, supporting either the fully coupled approach used here \cite{mccabe_development_2026} or a future decoupled formulation should the design space grow too large for full coupling.
\fi

%%%%%%%%%%%%%%%%%%%%%%%
\subsubsection{Parameter Suitability}
\ifdefined\DISSERTATION
    The local parameter sensitivity analysis of \Cref{remdo:sec:sensitivities-local} already yields significant insight into the choice of parameters.
    Interpreting those results more broadly, parameters with high design sensitivity $\left(\frac{dx^*}{dp}\right)$, such as dimensional aspect ratios, wave conditions, and material properties, are especially important for engineers to pin down early, since they affect the optimal design; where relevant, designers should consider holding their nondimensional forms constant across prototype scales, extending the common practice of Froude scaling for wave conditions.
    Parameters with high objective sensitivity $\left(\frac{dJ^*}{dp}\right)$, such as operational wave conditions, PTO efficiency, and financing rates in LCOE minimization, are most important for funders to determine early and for research agencies to target through R\&D incentives, since they affect economic viability.
    Notably, other parameters, such as PTO price, storm wave conditions, and dimensional variables, are more relevant to design cost minimization, highlighting that incentives and innovation strategies for Powering the Blue Economy applications may differ from those for utility-scale energy.
    The larger number of highly sensitive parameters in cost minimization further implies more diverse avenues for impactful innovation in PBE applications than in utility-scale energy.
    
    Further useful analysis includes a sensitivity analysis of hyperparameters, following \Cref{remdo:sec:sensitivities-local}, to ensure numerical convergence and adequate tuning, and a global parameter sensitivity analysis accounting for both nonlinearities across operating points and interaction effects among uncertain parameters.
\else
    Interpreting the parameter sensitivities (\Cref{remdo:sec:sensitivities-local}) broadly, parameters with high design sensitivity (aspect ratios, wave conditions, material properties) are most important for engineers to fix early,
    which should ideally be held constant in nondimensional form across scales,
    while those with high objective sensitivity (wave conditions, PTO efficiency, financing rates) are most important for funders and R\&D incentives.
    The differing high-sensitivity parameters for cost minimization, and their greater number, suggest that 
    innovation strategies for Powering the Blue Economy applications should differ from, and offer more avenues than, those for utility-scale energy.
\fi

%%%%%%%%%%%%%%%%%%%%%%%
\subsubsection{Constraint and Requirement Suitability}
\ifdefined\DISSERTATION
    By including \resultsRE[numLinConstraints]~linear and \resultsRE[numNonlinConstraints]~nonlinear constraints, this optimization captures more realistic design, model, and operational requirements than any other WEC optimization study known to the authors.
    The active constraints in \Cref{remdo:fig:constraint-activity} fall broadly into survivability: structural (design-variable lower bounds and factors of safety in the float, column, and damping plate) and amplitude-based (float rising above spar, linear theory, and float slamming).
    Constraints related to ballast and pitch stability are inactive, confirming these are not primary innovation avenues.
    While survivability is a known critical WEC requirement, it is rarely incorporated into optimization studies; this study demonstrates and quantifies its effect.
    For example, the Lagrange multipliers of \Cref{remdo:fig:lagrange-multiplier} suggest that improving the damping plate's structural design is worth about 5 times more than equivalent column improvements and 10 times more than float improvements.
    
    Further survivability innovation is warranted, both incremental and transformative.
    Incremental improvements include more structurally and hydrodynamically performant detail designs that obey these constraints while reducing material and sustaining energy production, for example, performing FEA or adding stiffeners to justify removing more material, replacing welded connections with pins to reduce moment transfer, or modifying the float--spar attachment geometry to allow more travel at the same bulk dimensions.
    Transformative improvements include creative concepts that avoid these constraints altogether, storm-survival strategies such as submerging the WEC, changing shape with control surfaces, or altering the PTO control strategy to reduce force and amplitude, whose viability can be readily estimated by weighing the additional costs (e.g.\ actuator and control hardware and software) against the performance gains, again using the Lagrange multipliers of \Cref{remdo:fig:lagrange-multiplier}.
    
    That said, many considerations remain uncaptured, requiring caution before asserting that industry should pursue the optimal design found here.
    Design variables are assumed continuous, whereas structural thicknesses and, to a lesser extent, generator specifications are constrained to discrete commercially available sizes.
    Additional requirement omissions for future work include mooring loads, structural resonant mode frequencies, and overturning moment in the storm design load case.
    Requirements such as installation and manufacturing feasibility, environmental and regulatory acceptability, reliability, and social benefit are harder to incorporate but nonetheless important;
    for instance, maximum dimensions may be limited by roadway transport, as for wind turbines, constraining component dimensions to 4.3--4.6~m \cite{cotrell_analysis_2014}, though relating maximum transported component size to maximum installed WEC dimensions requires manufacturability and on-site assembly analysis beyond this scope.
    For guidance on requirements and constraints in WEC design optimization, the reader is referred to \cite{bull_systems_2017,babarit_stakeholder_2017} for a comprehensive breakdown and mapping of WEC requirements, and to \cite{trueworthy_wave_2020} for how developers prioritize requirements in practice.
\else
    With \resultsRE[numLinConstraints]~linear and \resultsRE[numNonlinConstraints]~nonlinear constraints, this optimization captures more realistic requirements than any other WEC optimization known to the authors.
    The active constraints (\Cref{remdo:fig:constraint-activity}) are all survivability-related, such as structural factors of safety and amplitude limits, while ballast and pitch-stability constraints are inactive.
    Although survivability is a critical but rarely optimized WEC requirement, this study quantifies its effect:
    the Lagrange multipliers (\Cref{remdo:fig:lagrange-multiplier}) indicate that improving the damping plate is worth about 5 and 10 times more than improving the column and float respectively,
    motivating both incremental detail-design improvements and transformative storm-survival concepts (submersion, control surfaces, or amplitude-reducing control), whose viability can be screened using the same multipliers.
    Many considerations remain uncaptured, however, such as discrete component sizing, mooring loads, structural resonant modes, storm overturning moment, and manufacturing, regulatory, and transport limits \cite{cotrell_analysis_2014}, so caution is warranted before pursuing this optimum; see \cite{bull_systems_2017,babarit_stakeholder_2017,trueworthy_wave_2020} for comprehensive WEC requirement guidance.
\fi

%%%%%%%%%%%%%%%%%%%%%%%
\subsubsection{Objective Suitability}
\ifdefined\DISSERTATION
    Due to LCOE's high uncertainty, this study performs multi-objective optimization on cost and power and includes local and limited global sensitivity analysis.
    While this overcomes many concerns, limitations remain: prices for steel, powertrain force, and powertrain power must still be assumed, which could be addressed by adding objectives for each cost separately, though moving beyond two objectives limits the intuition and visualization of results, and the authors believe the two chosen represent a good balance.
    Alternatively, a global parameter sensitivity analysis of the price parameters, paired with uncertainty quantification and robustness analysis, would complete the picture.
    Although the cost formulation is uncertain, it still represents a significant improvement over the cost proxies typically used in WEC optimization, such as wetted surface area and hull volume.
    Future work running the MDOcean optimization with those simplified metrics would be useful for comparison and to assess the suitability of various proxy metrics.
\else
    Because of LCOE's uncertainty, this study optimizes cost and power as separate objectives with supporting sensitivity analysis.
    Prices for steel and powertrain force and power must still be assumed, which are addressable via more objectives or a global price-parameter sensitivity analysis with uncertainty quantification,
    but the two-objective formulation balances rigor against the interpretability lost beyond two objectives.
    The cost model, though uncertain, still improves substantially on the surface-area and hull-volume proxies common in WEC optimization, against which it would be useful to compare in future work.
\fi
% 	\paragraph{Not including powertrain impedance/generator model/real control co-design

% 	\paragraph{Grid scale does not include PBE}

%\paragraph{Does this design have high TPL? What are the risks and uncertainties?}
%use \url{https://tpl.nrel.gov/assessment} and \url{https://www.osti.gov/biblio/1825239} for RM3 TPL.

\subsection{Future Work}
\ifdefined\DISSERTATION
    Future improvements span two axes: enhancements to the model and enhancements to the optimization formulation, each applicable to either the existing setup or a new one.
    \Cref{remdo:tab:future-studies} organizes the most promising directions along these axes.

%% table contents helpers
\newcommand{\speedupsRE}{
    Examples of changes that would make it even faster:
    1) use gradient sparsity pattern to get rid of FD evaluations, or get rid of FD altogether by using matlab automatic differentiation
    2) make multi body code simpler more like single body
}

\newcommand{\achievableDesignStudiesRE}{
    Example design studies that could be done with the model more-or-less as it is now:

	1) more rigorously determine the effect of subsystem coupling - compare to sequential optimization of hydro then PTO sizing then structural - would let you know the benefit of coupled optimization for each subsystem, if it is large enough to justify the extra computational cost

	2) survivability - compare PTO-free vs PTO-braked in storm, to answer question of which type of survival mode is preferable; excitation loads using MEEM when float is sunk (for now could approximate as float with larger draft, once weird-region MEEM exists use that); allow decision on which sea states are operational vs survival could answer the question of whether it makes economic sense to go into survival mode more often.

	3) more detailed PTO study: include a drivetrain with linear dynamics reflecting spring, flywheel, gear ratio, and related effects. -  If you assume the flywheel inertia and spring stiffness are not controllable per sea state, would be useful as-is for finding the best inertia/stiffness across sea states to pair with the controller that can adjust per sea state.
    But if you assume adjustable drivetrain, as has been achieved with magnetic spring (cite), this would either need to introduce constraints or a cost term related to these drivetrain parameters, since otherwise the solution should go to zero damping and stiffness equal to what the unconstrained optimal control stiffness is.
    Could answer the interesting question of how much an adjustable PTO is worth over a non-adjustable one.
    Could also provide a preliminary answer to the question of optimum gear ratio (low GR means high generator torque so more force saturation for a given generator, but less friction.
    Vs high gear ratio has less force saturation but higher friction), although a generator model would be needed to more conclusively answer the question (discussed below).

    4) Add $N_{WEC}$ as a design variable and see when it's preferable to do many smaller WECs vs one bigger WEC

    5) How do my results compare to optimizing for easier cost proxies (ie structural cost as surface area vs actual force-informed thicknesses)
}

\newcommand{\modelTrustBuildersRE}{
    Examples of model enhancements that would not really unlock new types of design studies on their own, but would make this and any other design study more accurate and trustworthy/realistic:

	1) getting force with different nonlinear method: would make storm force more accurate, important since that is probably the least-theortically-defensible assumption of MDOcean

	2) more structures features: could use fatigue lifetime and in econ have lifetime = min(p.lifetime, fatigue lifetime) to be less conservative in structural sizing (currently uses endurance limit, which implies infinite fatigue lifetime)

    3) drag integral

    4) static friction describing function - ability to model static friction in drivetrain, and hydraulic check valves

    5) use extreme statistics to assess amplitude constraints, instead of max nonzero jpd
}

\newcommand{\modelStudyEnablersRE}{
    Example model additions and the types of design studies they would unlock:

	1) Multibody: would enable explicit optimization of spar dimensions, which is important because the minimum damping plate diameter constraint is frequently active.

	2) Irregular waves: would enable more realistic fatigue estimation (i.e.\ Miner's rule) and power-variability assessment (i.e.\ with energy storage), allowing comparison of variability reduction from design changes versus added batteries.
This would also provide more complete results for PTO impedance-matching bandwidth, including whether some PTO types or WEC shapes yield better $Z(\omega)$ matching with each other or with the sea state.

	3) model different WEC archetypes: would allow consistent comparison of WEC archetypes in early concept design and start to address the design convergence problem.
Would require not only MEEM for the different WEC types, but also semi-analytical structural models for each one, or the decision to integrate a lightweight FEA with plate/shell element capabilities (ie pynite).

	4) make lifetime a design variable, and adjust storm-wave size accordingly. This could address what OPEX cost threshold would justify intentionally designing ``single-use'' WECs, which might be considered in PBE applications.
Would require integration with WDRT or similar to get the environmental contour for a given lifetime.

	5) multi region MEEM: could not only make hydro coeffs more accurate representing the truncated cone shape, but would allow qualitatively different hydro coeff vs frequency shapes, which would unlock different

	6) generator physics and model - building off the PTO CCD study that is easily realizable with the current model, a generator model would unlock control co-design with the generator itself, for example whether a high-torque (expensive) generator is worth it compared to a low torque generator.
    The torque limit would essentially impose a constraint on the $F_{\text{max}}$ design var, the core losses introduce nonlinear damping, and impose a relation between generator torque limit and max/min generator inertia.
    A simplified generator model is explored in the RM3 CCD study \cite{anderson_re-imagining_2024}, or a full generator model as in the offshore wind MDO study \cite{barter_beyond_2023}.

    7) structures constraint that accounts for other types of materials (ie brittle fracture (Mohr's circle) for concrete, however inflatables are modeled, composites, inflatable polyurethane coated nylon fabric, which \cite{roberts_bringing_2021} found to be more optimal than steel/reinforced concrete/fiberglass/rubber).
    My relations between force and stress should still be good, but the relation between stress and FOS (the material limit state) is what needs to change.
}

% \begin{longtable}{
%     P{0.1\linewidth}|
%     P{0.3\linewidth}
%     P{0.6\linewidth}}
%                             &  Existing optimization formulation & New optimization formulation \\ \hline
%      Existing model features& \speedupsRE                        & \achievableDesignStudiesRE     \\
%      New model features     & \modelTrustBuildersRE              & \modelStudyEnablersRE        \\
%     \caption{Future model and optimization improvements}
%     \label{remdo:tab:future-studies}
% \end{longtable}
\else
    Future improvements span two axes, model fidelity and optimization formulation, each applicable to the existing setup or a new one.
    With the current model, the highest-value studies are a sequential-versus-coupled comparison to quantify the benefit of full coupling per subsystem, a survivability study comparing storm survival strategies, a detailed PTO drivetrain and gear-ratio study, an array-sizing study adding $N_{WEC}$ as a design variable, and a comparison against simplified cost proxies.
    Accuracy-enhancing model additions include a more defensible nonlinear storm-force method, fatigue-lifetime-based structural sizing, and extreme-statistics amplitude assessment.
    The most transformative additions, such as multibody dynamics (enabling true spar optimization), irregular waves (enabling realistic fatigue and power-variability studies), additional WEC archetypes (addressing design convergence), multi-region MEEM hydrodynamics, and a generator physics model (enabling generator control co-design \cite{anderson_re-imagining_2024,barter_beyond_2023}), would unlock qualitatively new design studies.
    The authors expect the speed of the semi-analytical models to permit full coupling even at substantially larger scale, including the ambitious target of comparing multiple hydrodynamic architectures with realistic powertrain models.
\fi
